\documentclass[11pt]{article}
\usepackage[a4paper,margin=1in]{geometry}
\usepackage{amsmath,amssymb,amsthm,mathtools,bm}
\usepackage{microtype}
\usepackage{enumitem}
\usepackage[hidelinks]{hyperref}
\usepackage{xcolor}

\newtheorem{theorem}{Theorem}[section]
\newtheorem{proposition}[theorem]{Proposition}
\newtheorem{lemma}[theorem]{Lemma}
\newtheorem{corollary}[theorem]{Corollary}
\newtheorem{claim}[theorem]{Claim}
\newtheorem{problem}[theorem]{Problem}
\newtheorem{definition}[theorem]{Definition}
\newtheorem{remark}[theorem]{Remark}

\newcommand{\one}{\mathbf 1}
\newcommand{\cQ}{\mathcal Q}
\newcommand{\cL}{\mathcal L}
\newcommand{\cP}{\mathcal P}
\newcommand{\cB}{\mathcal B}
\newcommand{\pos}[1]{\left(#1\right)_{+}}
\newcommand{\ip}[2]{\left\langle #1,#2\right\rangle}
\newcommand{\Tr}{\operatorname{Tr}}

\title{The Region-II Odd-Cycle Project:\\A Self-Contained Handoff Note for the Range $1/2<p<2/3$}
\author{Prepared for continuation by Claude and collaborators}
\date{\today}

\begin{document}
\maketitle

\begin{abstract}
This note records the work done on the remaining range
\[
        \frac12<p<\frac23
\]
for the conjectured odd-cycle lower bound
\[
        t(C_m,W)\ge p^m-p(1-p)^{m-1},
        \qquad m\ge3\text{ odd},
\]
where $W$ is a graphon of edge density $p$.  The high-density range
$p\ge2/3$ and all rank-two-complement, in particular bipodal, graphons
are handled in the companion notes.  The present document focuses only
on the unresolved Region II, where $q=1-p\in(1/3,1/2)$.

The main outcome is not a complete proof of Region II.  Rather, the
problem is reduced, rigorously, to an explicit two-variable scalar
payment inequality involving the unique complement-compression eigenvalue
above $q$.  The reduction keeps the two compensation channels that were
lost in earlier path-certificate and total-coupling relaxations: direct
coupling to the frontier eigenvector and indirect coupling through safe
spectral directions.  These two channels collapse to an exact Huber-type
envelope.  The note also records endpoint asymptotics, compact-interior
large-cycle reduction, and the failed proof routes that should not be
reused without repair.

The clean final target is Problem~\ref{prob:final-huber}: prove
\[
        R_m\le C_m\psi(\xi,\rho)
\]
for all admissible $(q,\alpha,m)$, with all quantities explicitly defined
below.  This target contains no graphon, no operator, no unknown spectral
measure, and no eigenfunction-shape variables.
\end{abstract}

\tableofcontents

\section{Executive summary and status}

Let $W:[0,1]^2\to[0,1]$ be a symmetric graphon and let
\[
        p=t(K_2,W)=\int W,\qquad q=1-p.
\]
The conjecture is
\begin{equation}\label{eq:conj}
        t(C_m,W)\ge p^m-pq^{m-1}
        \qquad (m\ge3\text{ odd}).
\end{equation}
For $p\le1/2$ the right-hand side is nonpositive, so the nontrivial
region is $p>1/2$.

The project currently has the following status.
\begin{enumerate}[label=(\roman*)]
\item The range $p\ge2/3$ has a draft proof by the two-sided
spectral-shift and diagonal-mixture method.
\item All graphons whose complement has rank two are handled for every
$p$ and every odd $m$.  In particular, all bipodal graphons are safe.
\item The only unresolved density range is
\[
        \frac12<p<\frac23,
        \qquad\text{or equivalently}\qquad
        \frac13<q<\frac12.
\]
\item In this range, the complement-compression operator has at most one
``frontier'' eigenvalue above $q$.  After the safe case is removed, the
remaining problem is a one-frontier spectral-payment problem.
\item The one-frontier problem has been reduced to the scalar Huber
inequality in Problem~\ref{prob:final-huber}.
\end{enumerate}

The strongest honest statement about Region II is therefore:
\[
\boxed{\text{Region II is not solved, but it has been reduced to an explicit scalar inequality.}}
\]

\section{Operator set-up and spectral-shift identities}

Put
\[
        U=1-W.
\]
The integral operator with kernel $U$ will be denoted by $T_U$.  On the
orthogonal decomposition
\[
        L^2([0,1])=\mathbb R\one\oplus\one^\perp
\]
write
\begin{equation}\label{eq:blockTU}
        T_U=
        \begin{pmatrix}
        q&g^*\\
        g&A
        \end{pmatrix}.
\end{equation}
Here
\[
        g=T_U\one-q\one\in\one^\perp,
        \qquad
        A=P_{\one^\perp}T_U P_{\one^\perp}.
\]
Since $0\le U\le1$, we have $U^2\le U$ pointwise, hence
\[
        \|T_U\|_{\mathrm{HS}}^2=\int U^2\le \int U=q.
\]
The block decomposition gives
\begin{equation}\label{eq:budget}
        q^2+2\|g\|_2^2+\Tr(A^2)\le q,
\end{equation}
or equivalently
\begin{equation}\label{eq:HS-budget}
        \Tr(A^2)+2\|g\|_2^2\le pq.
\end{equation}

The operator $T_W=J-T_U$, where $Jf=\ip{f}{\one}\one$.  If we conjugate
by the sign flip on $\one^\perp$, the trace powers of $T_W$ are the trace
powers of
\begin{equation}\label{eq:MW}
        M_W=
        \begin{pmatrix}
        p&g^*\\
        g&-A
        \end{pmatrix}.
\end{equation}
The complement side is
\begin{equation}\label{eq:MU}
        M_U=
        \begin{pmatrix}
        q&g^*\\
        g&A
        \end{pmatrix}=T_U.
\end{equation}

For finite-rank step graphons, Schur's determinant formula gives
\[
 \det(I-zM_W)=\det(I+zA)\left(1-pz-z^2\ip{g}{(I+zA)^{-1}g}\right),
\]
\[
 \det(I-zM_U)=\det(I-zA)\left(1-qz-z^2\ip{g}{(I-zA)^{-1}g}\right).
\]
Expanding $-\log\det(I-zT)=\sum_{r\ge1}z^r\Tr(T^r)/r$ gives the following.

\begin{proposition}[One-sided spectral shift]\label{prop:one-sided-shift}
For every odd $m\ge3$,
\begin{equation}\label{eq:one-sided-shift}
        t(C_m,W)=p^m-\Tr(A^m)+m[z^m]\cL_W(z),
\end{equation}
where
\[
        \cL_W(z)=
        -\log\left(
        1-\frac{z^2\ip{g}{(I+zA)^{-1}g}}{1-pz}
        \right).
\]
The corresponding complement formula is
\[
        t(C_m,U)=q^m+\Tr(A^m)+m[z^m]\cL_U(z),
\]
with
\[
        \cL_U(z)=
        -\log\left(
        1-\frac{z^2\ip{g}{(I-zA)^{-1}g}}{1-qz}
        \right).
\]
Consequently
\begin{equation}\label{eq:two-sided-shift}
        t(C_m,W)+t(C_m,U)=p^m+q^m+S_m,
\end{equation}
where
\[
        S_m=m[z^m](\cL_W(z)+\cL_U(z)).
\]
\end{proposition}

\begin{proof}
The proof above is written for finite-rank step graphons.  General
graphons follow by approximating $U$ in $L^2$ by finite-rank step kernels.
Cycle densities and the finitely many coefficients used above are
continuous under Hilbert-Schmidt convergence.  The oddness of $m$ is used
only in the cancellation
\[
        -\Tr(A^m)+\Tr(A^m)=0
\]
in the two-sided identity.
\end{proof}

The scripts included with this package audit these identities on rational
block graphons and reproduce the expansion checks used in the earlier
notes.

\section{Why Region II has only one frontier eigenvalue}

Assume from now on
\[
        \frac13<q<\frac12.
\]
The case $\lambda_{\max}(A)\le q$ is already safe in the spectral-shift
framework: the positive spectrum of $A$ is then too small to overcome the
baseline term $pq^{m-1}$, and the shift coefficients are nonnegative.
Hence we consider the remaining case where $A$ has an eigenvalue above
$q$.

\begin{lemma}[One-frontier property]\label{lem:one-frontier}
In Region II, $A$ has at most one eigenvalue greater than $q$.
\end{lemma}

\begin{proof}
If $A$ had two eigenvalues $>q$, then
\[
        \Tr(A^2)>2q^2.
\]
But by \eqref{eq:HS-budget}, $\Tr(A^2)\le pq=q(1-q)$.  Since
$q>1/3$, we have $2q^2>q(1-q)$, a contradiction.
\end{proof}

Thus, after replacing $A$ by its spectral decomposition, we may assume
there is a unique eigenpair
\begin{equation}\label{eq:frontier-eigenpair}
        A\phi=\alpha\phi,
        \qquad q<\alpha<\frac12,
        \qquad \phi\perp\one,
        \qquad \|\phi\|_2=1.
\end{equation}
Every other eigenvalue is at most $q$ and, more importantly for the
one-frontier reduction, lies in the interval $[-L,L]$, where
\begin{equation}\label{eq:def-L}
        L=\sqrt{pq-\alpha^2}.
\end{equation}
Indeed the square-mass outside the frontier is at most $pq-\alpha^2$.
Since $\alpha>q>1/3$, we have
\[
        L^2=pq-\alpha^2<pq-q^2=q(1-2q)<q^2,
\]
so
\begin{equation}\label{eq:L-order}
        0\le L<q<\alpha<p.
\end{equation}

\section{A universal forced-variance inequality}

The following lemma explains why the frontier cannot float above $q$
without forcing degree fluctuation.

\begin{lemma}[Forced variance]\label{lem:forced-variance}
For a frontier eigenpair \eqref{eq:frontier-eigenpair},
\begin{equation}\label{eq:forced-variance}
        \|g\|_2^2\ge
        \frac{\alpha(\alpha-q)^2}{2(1-2\alpha)}.
\end{equation}
Consequently,
\begin{equation}\label{eq:frontier-upper}
        \alpha\le r(q):=\frac{\sqrt{q^2+4q}-q}{2}.
\end{equation}
\end{lemma}

\begin{proof}
Let
\[
        a=\int |\phi|,
        \qquad f=|\phi|.
\]
Write $\phi=\phi^+-\phi^-$ with $\phi^+,\phi^-\ge0$.  Since $U\ge0$,
\[
        \alpha=\ip{\phi}{T_U\phi}
        \le \ip{\phi^+}{T_U\phi^+}+\ip{\phi^-}{T_U\phi^-}.
\]
Since $U\le1$,
\[
        \ip{\phi^+}{T_U\phi^+}\le\left(\int\phi^+\right)^2,
        \qquad
        \ip{\phi^-}{T_U\phi^-}\le\left(\int\phi^-\right)^2.
\]
Because $\int\phi=0$, both integrals are $a/2$, and therefore
\begin{equation}\label{eq:a2ge2alpha}
        a^2\ge2\alpha.
\end{equation}

Also $\ip{f}{T_Uf}\ge\ip{\phi}{T_U\phi}=\alpha$, since
$4\ip{\phi^+}{T_U\phi^-}=\ip{f}{T_Uf}-\ip{\phi}{T_U\phi}\ge0$.
Write $f=a\one+h$ with $h\perp\one$.  Since $\alpha$ is the top
eigenvalue of $A$,
\[
        \ip{h}{Ah}\le\alpha\|h\|_2^2=\alpha(1-a^2).
\]
Using the block form \eqref{eq:blockTU},
\[
        \alpha\le qa^2+2a\ip{g}{h}+\ip{h}{Ah}
        \le qa^2+2a\|g\|_2\sqrt{1-a^2}+\alpha(1-a^2).
\]
Thus
\[
        (\alpha-q)a^2\le2a\|g\|_2\sqrt{1-a^2}.
\]
After squaring and using \eqref{eq:a2ge2alpha},
\[
        \|g\|_2^2\ge
        \frac{(\alpha-q)^2a^2}{4(1-a^2)}
        \ge
        \frac{\alpha(\alpha-q)^2}{2(1-2\alpha)}.
\]
This proves \eqref{eq:forced-variance}.

Combining \eqref{eq:forced-variance} with
$\alpha^2+2\|g\|_2^2\le pq$ gives
\[
        \alpha^2+\frac{\alpha(\alpha-q)^2}{1-2\alpha}\le q(1-q).
\]
Multiplying by $1-2\alpha>0$ and simplifying gives
\[
        (1-\alpha-q)(\alpha^2+q\alpha-q)\le0.
\]
Since $\alpha+q<1$, we obtain $\alpha^2+q\alpha-q\le0$, which is exactly
\eqref{eq:frontier-upper}.
\end{proof}

\section{The one-frontier defect estimate}

Decompose the degree fluctuation as
\begin{equation}\label{eq:g-decomp}
        g=c\phi+g_s,
        \qquad c=\ip{g}{\phi},
        \qquad g_s\perp\phi.
\end{equation}
For odd $m$, define
\begin{equation}\label{eq:kdef}
        k_m(\lambda)=\frac{p^{m-1}-\lambda^{m-1}}{p+\lambda}.
\end{equation}
The function $k_m$ is positive and decreasing on $[0,p)$.
Set
\begin{equation}\label{eq:AB-def}
        A_m=2L^{m-2}+m k_m(\alpha),
        \qquad
        B_m=2L^{m-2}+m k_m(L).
\end{equation}
Then $0<A_m<B_m$.

\begin{proposition}[One-frontier master defect]\label{prop:master-defect}
Let
\begin{equation}\label{eq:R-def}
        R_m=\alpha^m+L^m-pq^{m-1}.
\end{equation}
Suppose we know a lower bound
\[
        \|g_s\|_2^2\ge \tau.
\]
Then
\begin{equation}\label{eq:defect-master-simple}
\begin{split}
        t(C_m,W)-\bigl(p^m-pq^{m-1}\bigr)
        \ge -R_m+A_m c^2+B_m\tau.
\end{split}
\end{equation}
\end{proposition}

\begin{proof}
The frontier eigenvalue contributes a possible negative term $-\alpha^m$
to the one-sided formula \eqref{eq:one-sided-shift}.  The remaining
positive spectrum has square mass at most
\[
        L^2-2c^2-2\|g_s\|_2^2
\]
by the Hilbert-Schmidt budget \eqref{eq:HS-budget}.  Since all nonfrontier
eigenvalues have absolute value at most $L$, their positive odd moment is
at most
\[
        L^{m-2}\bigl(L^2-2c^2-2\|g_s\|_2^2\bigr).
\]
Thus the unshifted positive spectral excess is at most
\[
        \alpha^m+L^m-2L^{m-2}(c^2+\|g_s\|_2^2).
\]

The spectral-shift series has nonnegative coefficients in Region II.
The term linear in each squared coupling gives at least
\[
        m k_m(\alpha)c^2+m k_m(L)\|g_s\|_2^2,
\]
because $k_m$ is decreasing on $[0,p)$ and every safe positive eigenvalue
is at most $L$.  Combining these contributions with the baseline
$pq^{m-1}$ gives \eqref{eq:defect-master-simple}.
\end{proof}

The remaining task is therefore to force enough payment
$A_mc^2+B_m\tau$ from the pointwise constraints $0\le U\le1$.

\section{Frontier-eigenfunction geometry and two coupling channels}

Define
\begin{equation}\label{eq:shape-defs}
        z=\left(\int|\phi|\right)^2,
        \qquad
        b=\ip{|\phi|-\sqrt z\one}{\phi},
\end{equation}
\[
        k=|\phi|-\sqrt z\one-b\phi,
        \qquad K=\|k\|_2^2.
\]
Replacing $\phi$ by $-\phi$ if necessary, assume
\begin{equation}\label{eq:bnonneg}
        b\ge0.
\end{equation}
Since $|\phi|$ has norm one and is decomposed orthogonally as
$\sqrt z\one+b\phi+k$, we have
\begin{equation}\label{eq:shape-budget}
        z+b^2+K=1.
\end{equation}
By \eqref{eq:a2ge2alpha}, $z\ge2\alpha$.  Put
\begin{equation}\label{eq:hdef}
        h=z-2\alpha,
        \qquad e=1-2\alpha.
\end{equation}
Then
\begin{equation}\label{eq:ebudget}
        h+b^2+K=e,
        \qquad h,b,K\ge0.
\end{equation}

\subsection{The direct channel}

\begin{lemma}[Direct endpoint]\label{lem:direct-endpoint}
The direct frontier coupling satisfies
\begin{equation}\label{eq:direct-upper}
        c\le u_c:=\frac{h-2\alpha b}{2\sqrt z}.
\end{equation}
There is also a lower endpoint, but it is not needed for the minimization
below.
\end{lemma}

\begin{proof}
Since $\phi^+=(|\phi|+\phi)/2\ge0$, and $U\le1$, we have pointwise
$T_U\phi\le T_1\phi^+=\int\phi^+=\sqrt z/2$.  Pairing with $\phi^+$ gives
\[
        \ip{\phi^+}{T_U\phi}\le \frac{\sqrt z}{2}\int\phi^+=\frac z4.
\]
On the other hand, because $T_U\phi=c\one+\alpha\phi$,
\[
        \ip{\phi^+}{T_U\phi}
        =\frac{c\sqrt z}{2}+\alpha\|\phi^+\|_2^2.
\]
Now
\[
        \|\phi^+\|_2^2=\frac{1+b}{2},
\]
since
$\|\phi^+\|_2^2-\|\phi^-\|_2^2=\ip{|\phi|}{\phi}=b$ and the two squares
sum to $1$.  Therefore
\[
        \frac{c\sqrt z}{2}+\frac{\alpha(1+b)}2\le\frac z4,
\]
which is \eqref{eq:direct-upper}.
\end{proof}

\subsection{The safe channel}

\begin{lemma}[Safe-channel constraint]\label{lem:safe-channel}
With
\begin{equation}\label{eq:H-def}
        H=\frac{(\alpha-q)z+(\alpha-L)K}{2\sqrt z},
\end{equation}
one has
\begin{equation}\label{eq:safe-constraint-main}
        bc+\ip{g_s}{k}\ge H.
\end{equation}
Consequently, when $K>0$,
\begin{equation}\label{eq:gs-lower}
        \|g_s\|_2^2\ge \frac{\pos{H-bc}^2}{K}.
\end{equation}
\end{lemma}

\begin{proof}
Since $U\ge0$,
\[
        4\ip{\phi^+}{T_U\phi^-}
        =\ip{|\phi|}{T_U|\phi|}-\ip{\phi}{T_U\phi}\ge0.
\]
Thus
\begin{equation}\label{eq:fUfgealpha}
        \ip{|\phi|}{T_U|\phi|}\ge\alpha.
\end{equation}
Using the decomposition $|\phi|=\sqrt z\one+b\phi+k$ and the block form
of $T_U$,
\[
\begin{split}
 \ip{|\phi|}{T_U|\phi|}
 &=qz+2\sqrt z\,(bc+\ip{g_s}{k})+\alpha b^2+\ip{k}{Ak}.
\end{split}
\]
The vector $k$ is orthogonal to both $\one$ and $\phi$, so it lies in the
safe spectral subspace, where $A\le L$ as a quadratic form.  Hence
$\ip{k}{Ak}\le LK$.  Combining this with \eqref{eq:fUfgealpha}, and using
$z+b^2+K=1$, gives
\[
        2\sqrt z\,(bc+\ip{g_s}{k})
        \ge (\alpha-q)z+(\alpha-L)K.
\]
This is \eqref{eq:safe-constraint-main}.  The bound
\eqref{eq:gs-lower} follows from Cauchy's inequality
$\ip{g_s}{k}\le\|g_s\|_2\sqrt K$.
\end{proof}

Combining Proposition~\ref{prop:master-defect} and Lemma~\ref{lem:safe-channel},
it is enough, for a fixed shape, to lower-bound the convex one-dimensional
payment
\begin{equation}\label{eq:Qexact}
        \cQ=
        \min_{c\le u_c}
        \left\{
        A_m c^2+\frac{B_m}{K}\pos{H-bc}^2
        \right\},
\end{equation}
with the usual limiting convention when $K=0$.

\section{Dual formulation and exact Huber elimination}

\subsection{Dual form for a fixed shape}

\begin{lemma}[Fixed-shape convex dual]\label{lem:fixed-dual}
For $K>0$, the minimum \eqref{eq:Qexact} equals
\begin{equation}\label{eq:fixed-dual}
        \max_{\lambda,\mu\ge0}
        \left\{
        \lambda H-\mu u_c
        -\frac{(\lambda b-\mu)^2}{4A_m}
        -\frac{\lambda^2K}{4B_m}
        \right\}.
\end{equation}
In particular,
\begin{equation}\label{eq:P0P1}
        \cQ\ge
        \max\left\{
        \frac{H^2}{b^2/A_m+K/B_m},
        \;B_m\frac{\pos{H-bu_c}^2}{K}
        \right\}.
\end{equation}
\end{lemma}

\begin{proof}
Use
\[
        \frac{B_m}{K}\pos{H-bc}^2
        =\max_{\lambda\ge0}
        \left\{\lambda(H-bc)-\frac{\lambda^2K}{4B_m}\right\}
\]
and impose $c\le u_c$ by the multiplier $\mu\ge0$, namely by adding
$\mu(u_c-c)$ in the minimization convention.  Minimizing the resulting
quadratic over $c$ gives \eqref{eq:fixed-dual} after a sign change in
$\mu$.

The first bound in \eqref{eq:P0P1} is obtained by taking $\mu=0$ and
optimizing in $\lambda$.  The second is obtained by setting
$\mu=\lambda b$ and optimizing in $\lambda$.
\end{proof}

\subsection{Huber elimination of all shape variables}

The dual bound above is useful, but still contains $b,K,h$.  The next
result removes them.

\begin{definition}\label{def:huber-params}
For Region II parameters $q,\alpha,m$, put
\[
        d=\alpha-q,
        \qquad f=\alpha-L,
        \qquad e=1-2\alpha,
\]
\begin{equation}\label{eq:C-xi-rho}
        C_m=\frac{B_m f\sqrt{2\alpha}\,e^2}{4\alpha^2},
        \qquad
        \xi=\frac{4\alpha^2 d}{e^2},
\end{equation}
\begin{equation}\label{eq:rho-def}
        \rho=\frac{A_m}{B_m}\frac{\sqrt\alpha}{2\sqrt2\,f},
\end{equation}
and define
\begin{equation}\label{eq:psi-def}
        \psi(\xi,\rho)=
        \min_{0\le v\le1}
        \left\{
        \rho v^2+\pos{\xi-v+v^2}
        \right\}.
\end{equation}
\end{definition}

\begin{theorem}[Exact Huber shape elimination]\label{thm:huber-elim}
For every graphon configuration in Region II,
\begin{equation}\label{eq:huber-payment}
        \cQ\ge C_m\psi(\xi,\rho).
\end{equation}
Consequently, the odd-cycle conjecture in Region II follows from the
scalar inequality
\begin{equation}\label{eq:scalar-huber}
        R_m\le C_m\psi(\xi,\rho)
\end{equation}
for every admissible $q,\alpha,m$.
\end{theorem}

\begin{proof}
We split according to the sign of $u_c$.

If $u_c\le0$, every admissible $c$ is nonpositive.  Since $b\ge0$,
$H-bc\ge H$, and therefore
\[
        \cQ\ge B_m\frac{H^2}{K}.
\]
From \eqref{eq:H-def},
\[
        \frac{H^2}{K}
        =\frac{((\alpha-q)z+(\alpha-L)K)^2}{4zK}
        \ge (\alpha-q)(\alpha-L)=df,
\]
because the inequality is equivalent to
$((\alpha-q)z-(\alpha-L)K)^2\ge0$.  Hence
\[
        \cQ\ge B_mdf.
\]
Since $\psi(\xi,\rho)\le\xi$ by taking $v=0$,
\[
        C_m\psi(\xi,\rho)
        \le C_m\xi
        =B_mf\sqrt{2\alpha}\,d
        \le B_mfd,
\]
so \eqref{eq:huber-payment} follows.

Now assume $u_c>0$.  Then any minimizer may be taken with $c\ge0$.
The inequality $c\le u_c$ gives
\[
        2c\sqrt z+2\alpha b\le h=e-b^2-K\le e-b^2.
\]
Using $z\ge2\alpha$, we get
\begin{equation}\label{eq:b-c-budget}
        b^2+2\alpha b+2c\sqrt{2\alpha}\le e.
\end{equation}
Thus
\[
        0\le c\le\frac{e}{2\sqrt{2\alpha}},
        \qquad
        b\le\beta(c):=
        \sqrt{\alpha^2+e-2c\sqrt{2\alpha}}-\alpha.
\]
Also, since $z\ge2\alpha$ and $z\le1$,
\[
        H=\frac{dz+fK}{2\sqrt z}
        \ge \frac{d\sqrt{2\alpha}}2+\frac f2K.
\]
Therefore
\[
        H-bc\ge
        \frac{d\sqrt{2\alpha}}2+\frac f2K-c\beta(c).
\]
For $w\ge0$,
\[
        \inf_{K>0}\frac{(w+fK/2)^2}{K}=2fw.
\]
Hence
\begin{equation}\label{eq:onevar-before-v}
        \cQ\ge
        \min_{0\le c\le e/(2\sqrt{2\alpha})}
        \left\{
        A_m c^2+2B_mf\pos{\frac{d\sqrt{2\alpha}}2-c\beta(c)}
        \right\}.
\end{equation}
Rationalizing $\beta(c)$ gives
\[
        \beta(c)=
        \frac{e-2c\sqrt{2\alpha}}
        {\sqrt{\alpha^2+e-2c\sqrt{2\alpha}}+\alpha}
        \le
        \frac{e-2c\sqrt{2\alpha}}{2\alpha}.
\]
Set
\[
        v=\frac{2c\sqrt{2\alpha}}e\in[0,1].
\]
Then
\[
        c\beta(c)
        \le
        \frac{e^2v(1-v)}{4\alpha\sqrt{2\alpha}},
\]
and
\[
        \frac{d\sqrt{2\alpha}}2-c\beta(c)
        \ge
        \frac{\sqrt{2\alpha}e^2}{8\alpha^2}
        (\xi-v+v^2).
\]
Also
\[
        A_mc^2=\frac{A_me^2}{8\alpha}v^2=C_m\rho v^2.
\]
Substituting into \eqref{eq:onevar-before-v} gives
\[
        \cQ\ge C_m
        \min_{0\le v\le1}\{\rho v^2+\pos{\xi-v+v^2}\},
\]
which is exactly \eqref{eq:huber-payment}.
The final assertion follows by combining this with
Proposition~\ref{prop:master-defect}.
\end{proof}

\section{Closed forms, inverse form, and dual certificates}

The Huber envelope is elementary.

\begin{proposition}[Closed form of $\psi$]\label{prop:psi-closed}
Let
\[
        \xi_c(\rho)=\frac{2\rho+1}{4(\rho+1)^2}.
\]
For $0\le\xi<1/4$, set
\[
        v_-(\xi)=\frac{1-\sqrt{1-4\xi}}2.
\]
Then
\begin{equation}\label{eq:psi-closed}
        \psi(\xi,\rho)=
        \begin{cases}
        \rho v_-(\xi)^2,
        &0\le\xi<\xi_c(\rho),\\[2mm]
        \displaystyle \xi-\frac1{4(1+\rho)},
        &\xi\ge\xi_c(\rho).
        \end{cases}
\end{equation}
The formula extends continuously at the transition.
\end{proposition}

\begin{proof}
On the set where $\xi-v+v^2\le0$, the objective equals $\rho v^2$ and is
minimized at the first root $v_-(\xi)$.  On the complementary set, the
objective is
\[
        (\rho+1)v^2-v+\xi,
\]
whose unconstrained minimizer is $v_*=1/(2(1+\rho))$ and whose value is
$\xi-1/(4(1+\rho))$.  The transition occurs when $v_*=v_-(\xi)$, which is
exactly $\xi=\xi_c(\rho)$.
\end{proof}

\begin{proposition}[Inverse form]\label{prop:inverse}
Let
\[
        T_c(\rho)=\frac{\rho}{4(1+\rho)^2}.
\]
For $T\ge0$, define
\[
        \Phi_\rho(T)=
        \begin{cases}
        \sqrt{T/\rho}-T/\rho,
        &0\le T\le T_c(\rho),\\[2mm]
        \displaystyle T+\frac1{4(1+\rho)},
        &T\ge T_c(\rho).
        \end{cases}
\]
Then
\[
        \psi(\xi,\rho)\ge T
        \quad\Longleftrightarrow\quad
        \xi\ge\Phi_\rho(T).
\]
Thus the final Region-II scalar target is equivalently
\begin{equation}\label{eq:inverse-final}
        \xi\ge \Phi_\rho(R_m/C_m).
\end{equation}
\end{proposition}

\begin{proof}
In the first branch, $\psi=\rho v_-^2$.  The condition
$\rho v_-^2\ge T$ is $v_-\ge\sqrt{T/\rho}$.  Since
$\xi=v_--v_-^2$ on this branch, the threshold is
$\sqrt{T/\rho}-T/\rho$.  In the second branch, the inequality is simply
$\xi-1/(4(1+\rho))\ge T$.
\end{proof}

There is also an exact dual representation.

\begin{proposition}[Dual representation of $\psi$]\label{prop:psi-dual}
For $\xi,\rho\ge0$,
\begin{equation}\label{eq:psi-dual}
        \psi(\xi,\rho)=
        \max_{0\le\lambda\le1}
        \left\{
        \lambda\xi-\frac{\lambda^2}{4(\rho+\lambda)}
        \right\}.
\end{equation}
\end{proposition}

\begin{proof}
Use the elementary identity
\[
        (t)_+=\max_{0\le\lambda\le1}\lambda t.
\]
Then
\[
\begin{split}
\psi(\xi,\rho)
&=\min_{v\in[0,1]}\max_{0\le\lambda\le1}
\{\rho v^2+\lambda(\xi-v+v^2)\}.
\end{split}
\]
The function is convex in $v$ and concave in $\lambda$, so minimax
interchange is valid.  For fixed $\lambda$, minimizing
$(\rho+\lambda)v^2-\lambda v+\lambda\xi$ over $v\in[0,1]$ gives the
unconstrained minimizer $v=\lambda/(2(\rho+\lambda))\in[0,1]$ and the value
in \eqref{eq:psi-dual}.
\end{proof}

This dual formula is the cleanest route to a computer-assisted completion:
on a parameter box, choose a rational $\lambda\in[0,1]$ and verify
\[
        C_m\left(\lambda\xi-\frac{\lambda^2}{4(\rho+\lambda)}\right)
        \ge R_m
\]
with interval arithmetic.

\section{Normalized scalar form}

It is useful to remove $p$ by putting
\begin{equation}\label{eq:normalized-vars}
        s=\frac qp,
        \qquad x=\frac\alpha p,
        \qquad y=\frac Lp.
\end{equation}
Then
\begin{equation}\label{eq:sxy-domain}
        s=x^2+y^2,
        \qquad
        \frac12<s<1,
        \qquad
        s<x<1,
        \qquad
        0\le y<s.
\end{equation}
The frontier upper condition is
\begin{equation}\label{eq:frontier-normalized}
        y^2\ge s(x-s).
\end{equation}
Conversely, these relations encode the Region-II one-frontier domain.

Define
\[
        E=1+s-2x=(1-x)^2+y^2,
\]
\[
        a_m=2y^{m-2}+m\frac{1-x^{m-1}}{1+x},
        \qquad
        b_m=2y^{m-2}+m\frac{1-y^{m-1}}{1+y},
\]
\[
        r_m=x^m+y^m-s^{m-1}.
\]
Then
\[
        R_m=p^m r_m,
        \qquad
        A_m=p^{m-2}a_m,
        \qquad
        B_m=p^{m-2}b_m.
\]
Moreover
\[
        \xi=\frac{4x^2(x-s)}{(1+s)E^2}.
\]
The exact Huber inequality can be written entirely in these variables by
substituting these expressions into \eqref{eq:scalar-huber}.  This is the
form used by the diagnostic scanner included in this package.

\section{Endpoint and large-cycle analysis}

The exact scalar inequality has been stress-tested analytically in the
singular regimes near $p=1/2$.  The point of this section is not to prove
the global scalar inequality, but to record the regimes already ruled out.

\subsection{Fixed odd length at $p\downarrow1/2$}

Let
\[
        e=1-2\alpha,
        \qquad d=\alpha-q=\kappa e,
\]
with $m$ fixed and $\kappa\to\kappa_0\in[0,1]$.  Then
\[
        \alpha=\frac{1-e}{2},
        \qquad
        q=\frac{1-e}{2}-\kappa e,
        \qquad
        p=\frac{1+e}{2}+\kappa e,
\]
and
\[
        L^2=eq-d^2
        =\frac e2-e^2\left(\frac12+\kappa+\kappa^2\right).
\]
Thus $L^m=o(e)$ for fixed odd $m$.

\begin{proposition}[Fixed-length endpoint]\label{prop:fixed-length}
For fixed odd $m$,
\[
        \frac{R_m}{p^me}\to 2((m-2)\kappa_0-1).
\]
If $\kappa_0>0$, then
\[
        \frac{C_m\psi(\xi,\rho)}{p^me}\to 2m\kappa_0.
\]
Hence the Huber payment dominates every positive limiting defect.
\end{proposition}

\begin{proof}
The expansion of $(\alpha/p)^m$ is
\[
        1-2m(1+\kappa_0)e+o(e),
\]
and the expansion of $(q/p)^{m-1}$ is
\[
        1-2(m-1)(1+2\kappa_0)e+o(e).
\]
This gives the expression for $R_m$.  If $\kappa_0>0$, then
$\xi\sim\kappa_0/e\to\infty$, and the closed form of $\psi$ gives
$\psi(\xi,\rho)=\xi+O(1)$.  Since $B_m\sim mp^{m-2}$ and
$C_m\xi=B_m(\alpha-L)\sqrt{2\alpha}\,d$, the payment limit is
$2m\kappa_0$.
\end{proof}

\subsection{Ultra-thin frontier scale}

The most delicate scale is
\begin{equation}\label{eq:ultra-thin-scale}
        d=a e^2,
        \qquad me\to T\in(0,\infty).
\end{equation}
Let
\[
        \rho_T=\frac{1-e^{-2T}}2
\]
and define
\[
        \Psi(a,\rho)=
        \min_{0\le C\le1/2}
        \left\{
        \rho C^2+\pos{2C^2-C+a/2}
        \right\}.
\]
The relation to $\psi$ is $\psi(a,\rho/2)=2\Psi(a,\rho)$ after $v=2C$.

\begin{proposition}[Ultra-thin limiting inequality]\label{prop:ultra-thin}
Under \eqref{eq:ultra-thin-scale}, after division by the common positive
scale $4p^me$, the limiting defect and payment are respectively
\[
        e^{-2T}(aT-1)
        \qquad\text{and}\qquad
        T\Psi(a,\rho_T).
\]
For all $a,T>0$,
\begin{equation}\label{eq:ultra-thin-ineq}
        T\Psi(a,\rho_T)
        \ge \frac12 e^{-2T}\pos{aT-1}.
\end{equation}
Thus the ultra-thin endpoint regime cannot generate a counterexample.
\end{proposition}

\begin{proof}
The asymptotic formulas follow from $L=O(\sqrt e)$,
$B_m\sim mp^{m-2}$,
\[
        \frac{A_m}{B_m}\to\frac{1-e^{-2T}}2,
        \qquad
        \xi\to a.
\]
It remains to prove \eqref{eq:ultra-thin-ineq}.

The closed form of $\Psi$ is obtained exactly as for $\psi$:
there is a transition
\[
        a_c(\rho)=\frac{\rho+1}{(\rho+2)^2},
\]
and the first branch satisfies $\Psi\ge\rho a^2/4$.  Put $z=aT>1$.
Since $(z-1)/z^2\le1/4$ and
$2\rho_T=1-e^{-2T}\ge T e^{-2T}$, the first branch gives
\[
        T\Psi\ge\frac{\rho_Tz^2}{4T}
        \ge\frac12e^{-2T}(z-1).
\]
In the second branch,
\[
        T\Psi=\frac z2-\frac{T}{4(\rho_T+2)}.
\]
The required inequality is
\[
        (1-e^{-2T})z+e^{-2T}\ge\frac{T}{2(\rho_T+2)}.
\]
For $T\le4$, the left side is at least $1$ and the right side is at most
$1$.  For $T>4$, using $z\ge a_c(\rho_T)T$ and the fact that
$\rho_T>(1-e^{-8})/2$ gives the same inequality by direct simplification.
\end{proof}

\subsection{Compact-interior large-cycle reduction}

\begin{proposition}[Large cycles away from endpoints]\label{prop:compact-tail}
For every $\eta>0$ there exists $M(\eta)$ such that the exact scalar target
\eqref{eq:scalar-huber} holds for all odd $m\ge M(\eta)$ whenever
\[
        \frac13+\eta\le q\le\frac12-\eta,
        \qquad
        q<\alpha\le r(q).
\]
\end{proposition}

\begin{proof}
On this compact set, both $\alpha/p$ and $L/p$ are bounded above by a
constant $\theta<1$.  Thus $R_m$ is exponentially small in $m$ except for
the affine displacement from the negative value at $\alpha=q$.  More
explicitly, $R_m(\alpha)$ is convex and
\[
        R_m'(\alpha)=m\alpha(\alpha^{m-2}-L^{m-2}).
\]
At $\alpha=q$, the value is negative by a uniform exponential amount.  On
the other hand, the small-$\xi$ branch of the Huber payment gives a
uniform quadratic lower bound
\[
        C_m\psi(\xi,\rho)\ge c_1mp^{m-2}(\alpha-q)^2
\]
for some $c_1=c_1(\eta)>0$.  The maximum of an affine function with slope
$O(mp^{m-1}\theta^{m})$ minus this quadratic term is smaller than the
negative base value for all sufficiently large $m$.  This proves the
claim.
\end{proof}

\section{Corrections and failed routes to avoid}

This section is included deliberately.  The Region-II project produced
several plausible but false shortcuts.  Future work should not reuse them
without repair.

\subsection{The shape-free total-coupling relaxation is too weak}

The forced-variance bound \eqref{eq:forced-variance} is true, but using
only $\|g\|_2^2$ loses the distinction between direct frontier coupling and
safe-channel coupling.  In the two-clique family, the naive rank-one-style
bound
\[
        \|g\|_2^2
        \gtrsim
        \frac{p\alpha(\alpha-q)^2}{(p-
        \alpha)^2}
\]
fails: the exact ratio is $9s+O(s^2)\to0$ along the relevant
parametrization.  The included checker reproduces this diagnostic exactly
with rational arithmetic.

\subsection{The smooth Huber lower bound is not globally strong enough}

The smooth lower bound
\[
        \psi(\xi,\rho)\ge
        \frac{(\sqrt{1+4\rho\xi}-1)^2}{4\rho}
\]
is valid, but the resulting sufficient inequality is false in Region II.
For example, with
\[
        m=13,
        \qquad q=0.498667598218,
        \qquad \alpha=0.499255849143,
\]
the exact Huber payment exceeds $R_m$, but the smooth payment does not.
Thus the exact piecewise Huber envelope is essential.

\subsection{A tempting recurrence for normalized powers is false}

In normalized variables, define
\[
        r_m=x^m+y^m-s^{m-1},
        \qquad s=x^2+y^2.
\]
A tempting but false recurrence is
\[
        r_{m+2}=sr_m-x^2y^2(x^{m-2}+y^{m-2}).
\]
The correct recurrence is
\begin{equation}\label{eq:correct-recurrence}
        r_{m+2}=sr_m-x^2y^2(x^{m-2}+y^{m-2})+s^m(1-s).
\end{equation}
The missing positive term $s^m(1-s)$ invalidates the earlier attempted
propagation proof for a direct-endpoint scalar inequality.

\subsection{Fixed-dual branch inequalities remain unproved}

The dual choice
\[
        \lambda_0=\min\{1,2\rho\xi\}
\]
reduces the exact Huber target to two explicit branch inequalities.  These
look numerically plausible in many tests, but we do not have a proof.  The
branch inequalities should be treated as diagnostics, not as established
lemmas.

\section{The final target for continuation}

The exact remaining scalar target is the following.

\begin{problem}[Final Huber payment inequality]\label{prob:final-huber}
Let $m\ge3$ be odd and let
\[
        \frac13<q<\frac12,
        \qquad
        q<\alpha\le\frac{\sqrt{q^2+4q}-q}{2}.
\]
Put $p=1-q$, $L=\sqrt{pq-\alpha^2}$, and define $A_m,B_m,R_m,C_m,\xi,\rho$
by \eqref{eq:AB-def}, \eqref{eq:R-def}, \eqref{eq:C-xi-rho}, and
\eqref{eq:rho-def}.  Prove
\[
        R_m\le C_m\psi(\xi,\rho).
\]
\end{problem}

A proof of Problem~\ref{prob:final-huber}, combined with the high-density
and rank-two-complement notes, would complete the original odd-cycle
conjecture.

The most realistic continuation strategies are:
\begin{enumerate}[label=(\alph*)]
\item Use the exact inverse form \eqref{eq:inverse-final} and prove the
piecewise scalar inequality directly.
\item Use the dual formula \eqref{eq:psi-dual}; on parameter boxes choose
rational $\lambda$ values and certify
\[
        C_m\left(\lambda\xi-\frac{\lambda^2}{4(\rho+\lambda)}\right)
        \ge R_m.
\]
\item Make Proposition~\ref{prop:compact-tail} quantitative, combine it
with explicit endpoint neighbourhoods from Propositions~\ref{prop:fixed-length}
and \ref{prop:ultra-thin}, and finish the remaining compact finite problem
with exact interval or Bernstein arithmetic.
\end{enumerate}

\appendix

\section{Program guide}

The package accompanying this note contains the following files.
\begin{itemize}
\item \texttt{regionII\_baton\_note.tex}: this note.
\item \texttt{regionII\_baton\_note.pdf}: compiled PDF.
\item \texttt{regionII\_handoff\_checker.py}: a lightweight checker for the
Region-II formulas, diagnostics, and numerical scalar target.
\item \texttt{clean\_two\_sided\_shift\_checker.py}: the earlier exact checker
for the two-sided spectral-shift identities, expansion formulae, mixture
identity, kernel form, and finite high-density certificates.
\item \texttt{README\_regionII\_baton.md}: instructions and current status.
\end{itemize}

The checker \texttt{regionII\_handoff\_checker.py} deliberately separates
exact algebraic checks from numerical stress tests.  The numerical scans
are not proof certificates.  Their purpose is to help locate the hard
transition regimes and to test any proposed analytic inequality before it
is written into a proof.

\section{A compact list of the principal formulae}

For quick reference:
\[
        p=1-q,
        \qquad
        L=\sqrt{pq-\alpha^2},
        \qquad
        d=\alpha-q,
        \qquad
        f=\alpha-L,
        \qquad
        e=1-2\alpha.
\]
\[
        R_m=\alpha^m+L^m-pq^{m-1}.
\]
\[
        A_m=2L^{m-2}+m\frac{p^{m-1}-\alpha^{m-1}}{p+\alpha},
        \qquad
        B_m=2L^{m-2}+m\frac{p^{m-1}-L^{m-1}}{p+L}.
\]
\[
        C_m=\frac{B_mf\sqrt{2\alpha}e^2}{4\alpha^2},
        \qquad
        \xi=\frac{4\alpha^2d}{e^2},
        \qquad
        \rho=\frac{A_m}{B_m}\frac{\sqrt\alpha}{2\sqrt2 f}.
\]
\[
        \psi(\xi,\rho)=
        \min_{0\le v\le1}\{\rho v^2+(\xi-v+v^2)_+\}.
\]
The final target is
\[
        R_m\le C_m\psi(\xi,\rho).
\]

\end{document}
